\section{附录B：编程基础知识速查}\label{ux9644ux5f55bux7f16ux7a0bux57faux7840ux77e5ux8bc6ux901fux67e5} 本附录提供了智慧水利平台开发中常用的编程语言基础知识速查，帮助读者快速回顾和掌握关键语法和概念。 \subsection{HTML5基础}\label{html5ux57faux7840} \subsubsection{文档结构}\label{ux6587ux6863ux7ed3ux6784} \begin{verbatim} <!DOCTYPE html> <html lang="zh-CN"> <head> <meta charset="UTF-8"> <meta name="viewport" content="width=device-width, initial-scale=1.0"> <title>文档标题</title> <link rel="stylesheet" href="styles.css"> <script src="script.js"></script> </head> <body> <!-- 页面内容 --> <h1>标题</h1> <p>段落</p> </body> </html> \end{verbatim} \subsubsection{常用标签}\label{ux5e38ux7528ux6807ux7b7e} \begin{longtable}[]{@{}ll@{}} \toprule\noalign{} 标签 & 描述 \\ \midrule\noalign{} \endhead \bottomrule\noalign{} \endlastfoot \texttt{\textless{}h1\textgreater{}} - \texttt{\textless{}h6\textgreater{}} & 标题 \\ \texttt{\textless{}p\textgreater{}} & 段落 \\ \texttt{\textless{}div\textgreater{}} & 块级容器 \\ \texttt{\textless{}span\textgreater{}} & 行内容器 \\ \texttt{\textless{}a\textgreater{}} & 超链接 \\ \texttt{\textless{}img\textgreater{}} & 图片 \\ \texttt{\textless{}ul\textgreater{}}, \texttt{\textless{}ol\textgreater{}}, \texttt{\textless{}li\textgreater{}} & 列表 \\ \texttt{\textless{}table\textgreater{}}, \texttt{\textless{}tr\textgreater{}}, \texttt{\textless{}td\textgreater{}} & 表格 \\ \texttt{\textless{}form\textgreater{}}, \texttt{\textless{}input\textgreater{}}, \texttt{\textless{}button\textgreater{}} & 表单元素 \\ \end{longtable} \subsection{CSS3基础}\label{css3ux57faux7840} \subsubsection{选择器}\label{ux9009ux62e9ux5668} \begin{verbatim} /* 元素选择器 */ div { color: blue; } /* 类选择器 */ .class-name { color: red; } /* ID选择器 */ #id-name { color: green; } /* 属性选择器 */ input[type="text"] { border: 1px solid gray; } /* 伪类选择器 */ a:hover { text-decoration: underline; } /* 子元素选择器 */ ul > li { list-style-type: square; } \end{verbatim} \subsubsection{盒模型}\label{ux76d2ux6a21ux578b} \begin{verbatim} div { width: 300px; height: 200px; padding: 20px; border: 1px solid black; margin: 30px; box-sizing: border-box; /* 包含padding和border */ } \end{verbatim} \subsection{JavaScript基础}\label{javascriptux57faux7840} \subsubsection{变量与数据类型}\label{ux53d8ux91cfux4e0eux6570ux636eux7c7bux578b} \begin{verbatim} // 变量声明 let name = "张三"; const age = 30; var isActive = true; // 数据类型 // 字符串 let str = "Hello"; // 数字 let num = 42; // 布尔值 let flag = false; // 数组 let arr = [1, 2, 3, 4]; // 对象 let obj = { name: "李四", age: 25 }; // 空值 let n = null; let u = undefined; \end{verbatim} \subsubsection{函数}\label{ux51fdux6570} \begin{verbatim} // 函数声明 function add(a, b) { return a + b; } // 箭头函数 const multiply = (a, b) => a * b; // 函数调用 let sum = add(5, 3); let product = multiply(4, 2); \end{verbatim} \subsubsection{条件与循环}\label{ux6761ux4ef6ux4e0eux5faaux73af} \begin{verbatim} // if条件 if (age > 18) { console.log("成年"); } else { console.log("未成年"); } // switch语句 switch (status) { case "online": console.log("在线"); break; case "offline": console.log("离线"); break; default: console.log("未知状态"); } // for循环 for (let i = 0; i < 5; i++) { console.log(i); } // while循环 let j = 0; while (j < 5) { console.log(j); j++; } // forEach遍历 arr.forEach(item => { console.log(item); }); \end{verbatim} \subsection{Python基础}\label{pythonux57faux7840} \subsubsection{变量与数据类型}\label{ux53d8ux91cfux4e0eux6570ux636eux7c7bux578b-1} \begin{verbatim} # 变量赋值 name = "张三" age = 30 is_active = True # 数据类型 # 字符串 text = "Hello World" # 数字 integer = 42 floating = 3.14 # 布尔值 flag = False # 列表 my_list = [1, 2, 3, 4, 5] # 元组 my_tuple = (1, 2, 3) # 字典 my_dict = {"name": "李四", "age": 25} # 集合 my_set = {1, 2, 3, 4} \end{verbatim} \subsubsection{函数}\label{ux51fdux6570-1} \begin{verbatim} # 函数定义 def greet(name): return f"Hello, {name}!" # 带默认参数的函数 def add(a, b=0): return a + b # 函数调用 message = greet("张三") result = add(5, 3) \end{verbatim} \subsubsection{条件与循环}\label{ux6761ux4ef6ux4e0eux5faaux73af-1} \begin{verbatim} # if条件 if age > 18: print("成年") elif age == 18: print("刚好成年") else: print("未成年") # for循环 for i in range(5): print(i) # while循环 j = 0 while j < 5: print(j) j += 1 # 列表推导式 squares = [x**2 for x in range(10)] \end{verbatim} \subsection{Java基础}\label{javaux57faux7840} \subsubsection{变量与数据类型}\label{ux53d8ux91cfux4e0eux6570ux636eux7c7bux578b-2} \begin{verbatim} // 基本数据类型 int number = 42; double pi = 3.14; boolean flag = true; char letter = 'A'; // 引用类型 String name = "张三"; int[] numbers = {1, 2, 3, 4, 5}; List<String> names = new ArrayList<>(); Map<String, Integer> ages = new HashMap<>(); \end{verbatim} \subsubsection{类与对象}\label{ux7c7bux4e0eux5bf9ux8c61} \begin{verbatim} // 类定义 public class Person { // 属性 private String name; private int age; // 构造函数 public Person(String name, int age) { this.name = name; this.age = age; } // 方法 public String getName() { return name; } public void setName(String name) { this.name = name; } public int getAge() { return age; } public void setAge(int age) { this.age = age; } } // 创建对象 Person person = new Person("张三", 30); \end{verbatim} \subsubsection{条件与循环}\label{ux6761ux4ef6ux4e0eux5faaux73af-2} \begin{verbatim} // if条件 if (age > 18) { System.out.println("成年"); } else { System.out.println("未成年"); } // switch语句 switch (status) { case "online": System.out.println("在线"); break; case "offline": System.out.println("离线"); break; default: System.out.println("未知状态"); } // for循环 for (int i = 0; i < 5; i++) { System.out.println(i); } // while循环 int j = 0; while (j < 5) { System.out.println(j); j++; } // 增强for循环 for (int num : numbers) { System.out.println(num); } \end{verbatim} \subsection{SQL基础}\label{sqlux57faux7840} \subsubsection{基本查询}\label{ux57faux672cux67e5ux8be2} \begin{verbatim} -- 查询所有列 SELECT * FROM users; -- 查询指定列 SELECT id, username, email FROM users; -- 条件查询 SELECT * FROM users WHERE age > 18; -- 排序 SELECT * FROM users ORDER BY created_at DESC; -- 分组 SELECT department, COUNT(*) FROM employees GROUP BY department; -- 连接查询 SELECT users.username, orders.order_number FROM users JOIN orders ON users.id = orders.user_id; \end{verbatim} \subsubsection{增删改}\label{ux589eux5220ux6539} \begin{verbatim} -- 插入数据 INSERT INTO users (username, email, age) VALUES ('张三', 'zhangsan@example.com', 25); -- 更新数据 UPDATE users SET age = 26 WHERE username = '张三'; -- 删除数据 DELETE FROM users WHERE id = 1; \end{verbatim}